%%%%%%%%%%%%%%%%%%%%%%%%%%%%%%%%%%%%%%%%%%%%%%%%%%%%%%%%%%%%
%% Lecture 14
%%%%%%%%%%%%%%%%%%%%%%%%%%%%%%%%%%%%%%%%%%%%%%%%%%%%%%%%%%%%

\lecture
{14}
{Friday, 25 September 2026}
{Introduction to Probability}
{Dr. Nishant Chandgotia}


%============================================================
% Weak Law of Large Numbers
%============================================================

\section{Weak Law of Large Numbers}

\subsection{Weak Law Under a Uniform Variance Bound}

Suppose that
\(
X_1,X_2,\ldots,X_n,\ldots
\)
are uncorrelated random variables satisfying
\(
\mathbb{E}(X_i)=\mu
\)
and
\(
\operatorname{Var}(X_i)\leq c
\)
for some constant $c>0$ and all $i$.

Let
\[
S_n=X_1+\cdots+X_n.
\]

Then
\(
\mathbb{E}(S_n)=n\mu
\)
and, since the random variables are uncorrelated,
\[
\operatorname{Var}(S_n)
=
\sum_{i=1}^n\operatorname{Var}(X_i)
\leq nc.
\]

Consequently,
\[
\begin{aligned}
\mathbb{E}
\left[
\left(
\frac{S_n}{n}-\mu
\right)^2
\right]
&=
\frac{1}{n^2}
\operatorname{Var}(S_n)\\
&\leq
\frac{c}{n}
\longrightarrow0.
\end{aligned}
\]

Hence
\[
\boxed{
\frac{S_n}{n}\longrightarrow\mu
\quad\text{in }L^2.
}
\]

Since convergence in $L^2$ implies convergence in probability,
we also have
\[
\boxed{
\frac{S_n}{n}\xrightarrow{P}\mu.
}
\]

\begin{observation}

The uniform variance bound
\(
\operatorname{Var}(X_i)\leq c
\)
is a strong assumption.

Our goal is to obtain a Weak Law of Large Numbers under weaker
conditions.

\end{observation}


%============================================================
% Truncation
%============================================================

\section{Truncation}

\subsection{Truncated Random Variables}

A standard way to remove the strong variance assumption is to
truncate the random variables.

For a sequence of random variables $X_1,X_2,\ldots$, define
\[
\overline{X}_i
=
X_i\mathbf{1}_{\{|X_i|\leq i\}}.
\]

More generally, for a truncation level $b_n>0$, define
\[
\overline{X}_{n,k}
=
X_{n,k}
\mathbf{1}_{\{|X_{n,k}|\leq b_n\}}.
\]

The idea is that the truncated variables have better controlled
second moments, while the probability that truncation actually
changes the variables can be made small.

We are interested in the relationship between
\(
\frac1n\sum_{i=1}^n\overline{X}_i
\)
and
\(
\frac1n\sum_{i=1}^nX_i.
\)

This leads naturally to the Weak Law for triangular arrays.


%============================================================
% Weak Law for Triangular Arrays
%============================================================

\section{Weak Law for Triangular Arrays}

\subsection{Statement}

\begin{theorem}{Weak Law for Triangular Arrays}
{weak-law-triangular-arrays}

For each $n\geq1$, let
\(
X_{n,k},
\qquad
1\leq k\leq n,
\)
be independent random variables.

Let
\(
b_n>0,
\qquad
b_n\longrightarrow\infty,
\)
and define the truncated variables
\[
\overline{X}_{n,k}
=
X_{n,k}
\mathbf{1}_{\{|X_{n,k}|\leq b_n\}}.
\]

Suppose that, as $n\to\infty$,

\begin{align}
\sum_{k=1}^n
\mathbb{P}(|X_{n,k}|>b_n)
&\longrightarrow0,
\tag{i}
\\
b_n^{-2}
\sum_{k=1}^n
\mathbb{E}\left[\overline{X}_{n,k}^{\,2}\right]
&\longrightarrow0.
\tag{ii}
\end{align}

Let
\[
S_n
=
X_{n,1}+\cdots+X_{n,n}
\]
and define
\[
a_n
=
\sum_{k=1}^n
\mathbb{E}\left[\overline{X}_{n,k}\right].
\]

Then
\[
\boxed{
\frac{S_n-a_n}{b_n}
\xrightarrow{P}0.
}
\]

\end{theorem}


\subsection{Proof}

\begin{proof}

Define
\[
\overline{S}_n
=
\overline{X}_{n,1}
+\cdots+
\overline{X}_{n,n}.
\]

We shall show that, for every $\varepsilon>0$,
\[
\mathbb{P}
\left(
\left|
\frac{S_n-a_n}{b_n}
\right|
>\varepsilon
\right)
\longrightarrow0.
\]

Since
\[
S_n-a_n
=
(S_n-\overline{S}_n)
+
(\overline{S}_n-a_n),
\]
we have
\[
\left\{
\left|
\frac{S_n-a_n}{b_n}
\right|
>\varepsilon
\right\}
\subseteq
\{S_n\neq\overline{S}_n\}
\cup
\left\{
\left|
\frac{\overline{S}_n-a_n}{b_n}
\right|
>\varepsilon
\right\}.
\]

Hence, by the union bound,
\begin{equation}
\mathbb{P}
\left(
\left|
\frac{S_n-a_n}{b_n}
\right|
>\varepsilon
\right)
\leq
\mathbb{P}(S_n\neq\overline{S}_n)
+
\mathbb{P}
\left(
\left|
\frac{\overline{S}_n-a_n}{b_n}
\right|
>\varepsilon
\right).
\tag{1}
\end{equation}

\medskip

\noindent
\textbf{Step 1: Probability that truncation changes the sum.}

If
\[
|X_{n,k}|\leq b_n,
\]
then
\[
\overline{X}_{n,k}
=
X_{n,k}
\mathbf{1}_{\{|X_{n,k}|\leq b_n\}}
=
X_{n,k}.
\]

Therefore,
\[
\{S_n\neq\overline{S}_n\}
\subseteq
\bigcup_{k=1}^n
\{\overline{X}_{n,k}\neq X_{n,k}\}.
\]

Moreover,
\[
\{\overline{X}_{n,k}\neq X_{n,k}\}
\subseteq
\{|X_{n,k}|>b_n\}.
\]

Hence
\[
\{S_n\neq\overline{S}_n\}
\subseteq
\bigcup_{k=1}^n
\{|X_{n,k}|>b_n\}.
\]

Consequently,
\[
\begin{aligned}
\mathbb{P}(S_n\neq\overline{S}_n)
&\leq
\mathbb{P}
\left(
\bigcup_{k=1}^n
\{|X_{n,k}|>b_n\}
\right)\\
&\leq
\sum_{k=1}^n
\mathbb{P}(|X_{n,k}|>b_n).
\end{aligned}
\]

By assumption~\textnormal{(i)},
\[
\boxed{
\mathbb{P}(S_n\neq\overline{S}_n)
\longrightarrow0.
}
\tag{2}
\]

\medskip

\noindent
\textbf{Step 2: Control of the truncated sum.}

By linearity of expectation,
\[
\begin{aligned}
\mathbb{E}(\overline{S}_n)
&=
\mathbb{E}
\left[
\sum_{k=1}^n\overline{X}_{n,k}
\right]\\
&=
\sum_{k=1}^n
\mathbb{E}(\overline{X}_{n,k})\\
&=
a_n.
\end{aligned}
\]

Thus
\[
\overline{S}_n-a_n
=
\overline{S}_n-\mathbb{E}(\overline{S}_n).
\]

Therefore, by Chebyshev's inequality,
\[
\begin{aligned}
\mathbb{P}
\left(
\left|
\frac{\overline{S}_n-a_n}{b_n}
\right|
>\varepsilon
\right)
&=
\mathbb{P}
\left(
|\overline{S}_n-\mathbb{E}(\overline{S}_n)|
>
\varepsilon b_n
\right)\\
&\leq
\frac{
\operatorname{Var}(\overline{S}_n)
}{
\varepsilon^2b_n^2
}.
\end{aligned}
\tag{3}
\]

Since the $X_{n,k}$ are independent, the truncated random variables
$\overline{X}_{n,k}$ are also independent. Hence
\[
\operatorname{Var}(\overline{S}_n)
=
\sum_{k=1}^n
\operatorname{Var}(\overline{X}_{n,k}).
\]

For any random variable $Y$,
\[
\operatorname{Var}(Y)
=
\mathbb{E}(Y^2)-(\mathbb{E}Y)^2
\leq
\mathbb{E}(Y^2).
\]

Therefore,
\[
\operatorname{Var}(\overline{X}_{n,k})
\leq
\mathbb{E}
\left[
\overline{X}_{n,k}^{\,2}
\right].
\]

Substituting this into (3), we obtain
\[
\begin{aligned}
\mathbb{P}
\left(
\left|
\frac{\overline{S}_n-a_n}{b_n}
\right|
>\varepsilon
\right)
&\leq
\frac{1}{\varepsilon^2b_n^2}
\sum_{k=1}^n
\operatorname{Var}(\overline{X}_{n,k})\\
&\leq
\frac{1}{\varepsilon^2b_n^2}
\sum_{k=1}^n
\mathbb{E}
\left[
\overline{X}_{n,k}^{\,2}
\right]\\
&=
\frac{1}{\varepsilon^2}
\left[
b_n^{-2}
\sum_{k=1}^n
\mathbb{E}
\left[
\overline{X}_{n,k}^{\,2}
\right]
\right].
\end{aligned}
\]

By assumption~\textnormal{(ii)},
\[
\boxed{
\mathbb{P}
\left(
\left|
\frac{\overline{S}_n-a_n}{b_n}
\right|
>\varepsilon
\right)
\longrightarrow0.
}
\tag{4}
\]

\medskip

\noindent
\textbf{Step 3: Conclusion.}

Combining (1), (2), and (4), we obtain
\[
\mathbb{P}
\left(
\left|
\frac{S_n-a_n}{b_n}
\right|
>\varepsilon
\right)
\longrightarrow0.
\]

Since this holds for every $\varepsilon>0$,
\[
\boxed{
\frac{S_n-a_n}{b_n}
\xrightarrow{P}0.
}
\]

Hence the proof is complete.

\end{proof}


%============================================================
% Weak Law for i.i.d. Random Variables
%============================================================

\section{A Weak Law Under a Tail Condition}

\subsection{Truncated Weak Law of Large Numbers}

\begin{theorem}{Weak Law of Large Numbers Under a Tail Condition}
{weak-law-tail-condition}

Let
\(
X_1,X_2,\ldots
\)
be i.i.d. random variables satisfying
\[
x\mathbb{P}(|X_1|>x)
\longrightarrow0
\qquad
\text{as }x\to\infty.
\]

Let
\(
S_n
=
X_1+\cdots+X_n
\)
and define
\(
\mu_n
=
\mathbb{E}
\left[
X_1
\mathbf{1}_{\{|X_1|\leq n\}}
\right].
\)

Then
\[
\boxed{
\frac{S_n}{n}-\mu_n
\xrightarrow{P}0.
}
\]

\end{theorem}


\subsection{Proof}

\begin{proof}

We apply Theorem Weak Law For Triangular Arryasto the triangular
array
\[
X_{n,k}=X_k,
\qquad
1\leq k\leq n,
\]
with
\[
b_n=n.
\]

Since
\[
X_1,X_2,\ldots
\]
are i.i.d., the variables
\[
X_{n,1},\ldots,X_{n,n}
\]
are independent for every $n$.

\medskip

\noindent
\textbf{Verification of condition (i).}

We have
\[
\begin{aligned}
\sum_{k=1}^n
\mathbb{P}(|X_{n,k}|>b_n)
&=
\sum_{k=1}^n
\mathbb{P}(|X_k|>n)\\
&=
n\mathbb{P}(|X_1|>n).
\end{aligned}
\]

By assumption,
\[
x\mathbb{P}(|X_1|>x)
\longrightarrow0.
\]

Taking $x=n$, we obtain
\[
n\mathbb{P}(|X_1|>n)
\longrightarrow0.
\]

Hence
\[
\boxed{
\sum_{k=1}^n
\mathbb{P}(|X_{n,k}|>b_n)
\longrightarrow0.
}
\tag{1}
\]

\medskip

\noindent
\textbf{Verification of condition (ii).}

We have
\[
\begin{aligned}
b_n^{-2}
\sum_{k=1}^n
\mathbb{E}
\left[
X_{n,k}^2
\mathbf{1}_{\{|X_{n,k}|\leq b_n\}}
\right]
&=
\frac{1}{n^2}
\sum_{k=1}^n
\mathbb{E}
\left[
X_k^2
\mathbf{1}_{\{|X_k|\leq n\}}
\right]\\
&=
\frac{1}{n}
\mathbb{E}
\left[
X_1^2
\mathbf{1}_{\{|X_1|\leq n\}}
\right].
\end{aligned}
\tag{2}
\]

It remains to show that
\[
\frac{1}{n}
\mathbb{E}
\left[
X_1^2
\mathbf{1}_{\{|X_1|\leq n\}}
\right]
\longrightarrow0.
\tag{3}
\]

Put
\[
Y=|X_1|.
\]

Then
\[
\mathbb{E}
\left[
X_1^2
\mathbf{1}_{\{|X_1|\leq n\}}
\right]
=
\mathbb{E}
\left[
Y^2
\mathbf{1}_{\{Y\leq n\}}
\right].
\]

Since
\[
Y^2\mathbf{1}_{\{Y\leq n\}}
\leq
Y^2\wedge n^2,
\]
we have
\[
\mathbb{E}
\left[
Y^2\mathbf{1}_{\{Y\leq n\}}
\right]
\leq
\mathbb{E}(Y^2\wedge n^2).
\]

Using the tail-integral formula,
\[
\mathbb{E}(Y^2\wedge n^2)
=
\int_0^{n^2}
\mathbb{P}(Y^2>t)\,dt.
\]

Making the substitution
\[
t=x^2,
\qquad
dt=2x\,dx,
\]
gives
\[
\mathbb{E}(Y^2\wedge n^2)
=
2\int_0^n
x\mathbb{P}(Y>x)\,dx.
\]

Consequently,
\[
\frac1n
\mathbb{E}
\left[
X_1^2
\mathbf{1}_{\{|X_1|\leq n\}}
\right]
\leq
\frac2n
\int_0^n
x\mathbb{P}(|X_1|>x)\,dx.
\tag{4}
\]

By assumption,
\[
x\mathbb{P}(|X_1|>x)
\longrightarrow0.
\]

Let $\varepsilon>0$. There exists $M>0$ such that
\[
x\mathbb{P}(|X_1|>x)
<
\varepsilon
\qquad
\text{for all }x\geq M.
\]

For $n>M$,
\[
\begin{aligned}
\frac1n
\int_0^n
x\mathbb{P}(|X_1|>x)\,dx
&=
\frac1n
\int_0^M
x\mathbb{P}(|X_1|>x)\,dx\\
&\quad+
\frac1n
\int_M^n
x\mathbb{P}(|X_1|>x)\,dx\\
&\leq
\frac1n
\int_0^M
x\mathbb{P}(|X_1|>x)\,dx\\
&\quad+
\frac{\varepsilon(n-M)}{n}.
\end{aligned}
\]

As $n\to\infty$, the first term tends to $0$, while
\[
\frac{\varepsilon(n-M)}{n}
\leq\varepsilon.
\]

Hence
\[
\limsup_{n\to\infty}
\frac1n
\int_0^n
x\mathbb{P}(|X_1|>x)\,dx
\leq\varepsilon.
\]

Since $\varepsilon>0$ is arbitrary,
\[
\frac1n
\int_0^n
x\mathbb{P}(|X_1|>x)\,dx
\longrightarrow0.
\]

It follows from (4) that
\[
\frac1n
\mathbb{E}
\left[
X_1^2
\mathbf{1}_{\{|X_1|\leq n\}}
\right]
\longrightarrow0.
\]

Thus condition (ii) is satisfied.

\medskip

Therefore, by Theorem~\ref{weak-law-triangular-arrays},
\[
\frac{S_n-a_n}{n}
\xrightarrow{P}0,
\]
where
\[
a_n
=
\sum_{k=1}^n
\mathbb{E}
\left[
X_k
\mathbf{1}_{\{|X_k|\leq n\}}
\right].
\]

Since the $X_k$ are identically distributed,
\[
\begin{aligned}
a_n
&=
n
\mathbb{E}
\left[
X_1
\mathbf{1}_{\{|X_1|\leq n\}}
\right]\\
&=
n\mu_n.
\end{aligned}
\]

Consequently,
\[
\frac{S_n-a_n}{n}
=
\frac{S_n-n\mu_n}{n}
=
\frac{S_n}{n}-\mu_n.
\]

Therefore,
\[
\boxed{
\frac{S_n}{n}-\mu_n
\xrightarrow{P}0.
}
\]

The proof is complete.

\end{proof}


\begin{remark}

The condition
\[
x\mathbb{P}(|X_1|>x)\longrightarrow0
\]
is the tail condition used above to obtain the weak law with the
truncated mean $\mu_n$.

\end{remark}


%============================================================
% Tail Conditions and Moments
%============================================================

\section{Tail Conditions and Moment Conditions}

\subsection{A Tail Condition Implies Lower Moments}

\begin{theorem}{Tail Condition and Fractional Moments}
{tail-condition-fractional-moments}

Let $X$ be a random variable.

If
\[
x\mathbb{P}(|X|>x)
\longrightarrow0
\qquad
\text{as }x\to\infty,
\]
then
\[
X\in L^{1-\varepsilon}
\qquad
\text{for every }0<\varepsilon<1.
\]

\end{theorem}


\begin{proof}

Suppose that
\[
x\mathbb{P}(|X|>x)
\longrightarrow0.
\]

We prove that
\[
X\in L^p
\qquad
\text{for every }0<p<1.
\]

Using the tail-integral formula, for $p>0$,
\[
\mathbb{E}|X|^p
=
p\int_0^\infty
x^{p-1}
\mathbb{P}(|X|>x)\,dx.
\]

Since
\[
x\mathbb{P}(|X|>x)
\longrightarrow0,
\]
there exists $M>0$ such that
\[
x\mathbb{P}(|X|>x)\leq1
\qquad
\text{for all }x\geq M.
\]

Therefore,
\[
\mathbb{P}(|X|>x)
\leq
\frac1x,
\qquad
x\geq M.
\]

Split the integral at $M$:
\[
\begin{aligned}
\mathbb{E}|X|^p
&=
p\int_0^M
x^{p-1}
\mathbb{P}(|X|>x)\,dx\\
&\quad+
p\int_M^\infty
x^{p-1}
\mathbb{P}(|X|>x)\,dx.
\end{aligned}
\]

For the first integral,
\[
\begin{aligned}
p\int_0^M
x^{p-1}
\mathbb{P}(|X|>x)\,dx
&\leq
p\int_0^M x^{p-1}\,dx\\
&=
M^p
<
\infty.
\end{aligned}
\]

For the second integral,
\[
\begin{aligned}
p\int_M^\infty
x^{p-1}
\mathbb{P}(|X|>x)\,dx
&\leq
p\int_M^\infty x^{p-2}\,dx\\
&=
\frac{p}{1-p}M^{p-1}\\
&<\infty,
\end{aligned}
\]
because
\[
0<p<1.
\]

Therefore,
\[
\mathbb{E}|X|^p<\infty.
\]

Hence
\[
X\in L^p
\qquad
\text{for every }0<p<1.
\]

Taking
\[
p=1-\varepsilon,
\]
we obtain
\[
\boxed{
X\in L^{1-\varepsilon}
\qquad
\text{for every }0<\varepsilon<1.
}
\]

\end{proof}


%============================================================
% Integrability Implies the Tail Condition
%============================================================

\subsection{Integrability Implies the Tail Condition}

\begin{theorem}{Integrability Implies Vanishing Tail}
{integrability-tail-condition}

Let $X$ be a random variable.

If
\[
X\in L^1,
\]
then
\[
\boxed{
x\mathbb{P}(|X|>x)
\longrightarrow0
\qquad
\text{as }x\to\infty.
}
\]

\end{theorem}


\begin{proof}

Suppose that
\[
X\in L^1.
\]

Then
\[
\mathbb{E}|X|<\infty.
\]

On the event
\[
\{|X|>x\},
\]
we have
\[
x<|X|.
\]

Therefore,
\[
x\mathbf{1}_{\{|X|>x\}}
\leq
|X|\mathbf{1}_{\{|X|>x\}}.
\]

Taking expectations,
\[
x\mathbb{P}(|X|>x)
\leq
\mathbb{E}
\left[
|X|\mathbf{1}_{\{|X|>x\}}
\right].
\]

For every $\omega$,
\[
|X(\omega)|
\mathbf{1}_{\{|X(\omega)|>x\}}
\longrightarrow0
\qquad
\text{as }x\to\infty.
\]

Moreover,
\[
0\leq
|X|\mathbf{1}_{\{|X|>x\}}
\leq
|X|,
\]
and
\[
\mathbb{E}|X|<\infty.
\]

Therefore, by the dominated convergence theorem,
\[
\mathbb{E}
\left[
|X|\mathbf{1}_{\{|X|>x\}}
\right]
\longrightarrow0.
\]

Consequently,
\[
0\leq
x\mathbb{P}(|X|>x)
\leq
\mathbb{E}
\left[
|X|\mathbf{1}_{\{|X|>x\}}
\right]
\longrightarrow0.
\]

Hence
\[
\boxed{
x\mathbb{P}(|X|>x)
\longrightarrow0.
}
\]

The proof is complete.

\end{proof}


%============================================================
% Relation Between the Two Conditions
%============================================================

\subsection{Relation Between Tail and Moment Conditions}

\begin{observation}

The two preceding results give the implications
\[
X\in L^1
\quad\Longrightarrow\quad
x\mathbb{P}(|X|>x)\longrightarrow0,
\]
while
\[
x\mathbb{P}(|X|>x)\longrightarrow0
\quad\Longrightarrow\quad
X\in L^p
\quad\text{for every }0<p<1.
\]

Thus the tail condition lies between integrability and the collection
of fractional moment conditions:
\[
L^1
\quad\Longrightarrow\quad
\left\{
x\mathbb{P}(|X|>x)\to0
\right\}
\quad\Longrightarrow\quad
\bigcap_{0<p<1}L^p.
\]

\end{observation}


%============================================================
% End of Lecture
%============================================================

\lectureend
